\chapter{Méthodologie}
\label{chapitre:2}
Ce chapitre présente le processus d'élaboration et de validation de modèles de classification de
données métagénomiques.

\section{Bases de données}
Les modèles sont entraînés et validés à partir de trois bases de données.

\begin{enumerate}
    \item RefSeq \cite{refseq} est une base de données de séquences biologiques de haute qualité
        et non redondante. Ces séquences sont utilisées pour générer des données synthétiques afin
        de simuler des échantillons métagénomiques réels.
    \item GTDB \cite{gtdb} est une base de données taxonomique de bactéries et archées basée sur
        une approche phylogénétique.
    \item NCBI Taxonomy \cite{ncbi_tax} est une base de donnée taxonomique pour tous les organismes
        étant représentés dans la base de données RefSeq. L'approche de cette base de donnée est
        moins rigoureuse que GTDB mais comprend davantage de données.
\end{enumerate}

\section{Encodage}
Les données tirées de RefSeq sont des fichiers FASTA qui contiennent des chaînes de caractères :
A, C, G, T, et N pour les nucléotides incertains. Puisque les réseaux neuronaux prennent en entrée
des matrices numériques, il est nécessaire de transformer ces données en un format acceptable.

\subsection{K-mer}
L'encodage en K-mer consiste à mesurer la fréquence de K-mers dans une séquence. On obtient un
dictionnaire qui associe à chaque K-mer une fréquence relative. Cette représentation entraîne une
perte de données en lien avec l'ordonnancement des nucléotides et ne retient que les fréquences
relative de sous-séquences. Il s'agit néanmoins d'une représentation efficace pour les modèles
simples.

\subsection{Jetonisation}
La jetonisation consiste à diviser une séquence en sous-séquences de longueurs identiques et à leur
assigner un identifiant numérique unique à chaque sous-séquence identique. La taille du vocabulaire
(c'est-à-dire, le nombre d'identifiants différents) équivaut à $4^J$, où $J$ est la longueur des
sous-séquences. Un désavantage de la jetonisation est qu'elle introduit des artefacts qui ne sont
pas présents dans la séquence originale, par exemple, en ajoutant des segments entre les
nucléotides.

\subsection{1 parmi N}
L'encodage 1 parmi N assigne un vecteur orthogonal à chaque type de nucléotides, comme le montre
l'équation \ref{eq:1_parmi_n}.

\begin{equation}
    \label{eq:1_parmi_n}
    \begin{split}
        A &= [1, 0, 0, 0] \\
        C &= [0, 1, 0, 0] \\
        G &= [0, 0, 1, 0] \\
        T &= [0, 0, 0, 1]
    \end{split}
\end{equation}

L'encodage 1 parmi N n'introduit pas d'artefact, contrairement à la jetonisation, mais il requière
davantage de ressources de calcul.

\section{Modèles}
Les modèles suivants sont validés.

\subsection{Modèle aléatoire}
Le modèle aléatoire consiste à assigner un taxon aléatoire pour chaque séquence. Ce modèle agit
comme référence pour les autres modèles.

\subsection{Forêt aléatoire}
La forêt aléatoire \cite{breiman2001} est un modèle d'apprentissage automatique classique
(c'est-à-dire non basé sur les réseaux neuronaux). Ce modèle utilise l'encodage de K-mers pour
classifier les séquences et agit comme une seconde référence pour comparer les performances des
modèles plus performants.

\subsection{Modèles de classification taxonomique classiques}
On utilise les modèles Kraken 2 \cite{kraken} et MMSeqs 2 \cite{mmseqs2} pour les comparer avec les
méthodes d'apprentissage automatique.

\subsection{Perceptron multicouche}

\subsection{Réseau neuronal convolutionnel}

\subsection{Autoencodeur débruiteur}

\subsection{Transformateur encodeur seul}
Le transformateur encodeur seul testé a une architecture définie par les hyperparamètres présentés
dans le tableau \ref{tableau:transformateur_hyperparametre}.

\begin{table}[hbt!]
    \centering
    \begin{tabular}{ll}
\hline
Nom de l'hyperparamètre & Valeur \\
\hline
Taille des jetons & $4$ \\
Taille du vocabulaire & $256$ \\
Nombre de couches & $4$ \\
Nombre de têtes & $4$ \\
Dimension de plongement & $128$ \\
\hline
    \end{tabular}
    \caption{Hyperparamètres du transformateur encodeur seul}
    \label{tableau:transformateur_hyperparametre}
\end{table}

\subsection{Mamba}
Le modèle Mamba testé a une architecture définie par les hyperparamètres présentés
dans le tableau \ref{tableau:mamba_hyperparametre}.

\begin{table}[hbt!]
    \centering
    \begin{tabular}{ll}
\hline
Nom de l'hyperparamètre & Valeur \\
\hline
Taille des jetons & $4$ \\
Taille du vocabulaire & $256$ \\
Nombre de couches & $4$ \\
Dimension de l'état & $16$ \\
Facteur d'expansion & $2$ \\
Taille de convolution & $4$ \\
\hline
    \end{tabular}
    \caption{Hyperparamètres du modèle Mamba}
    \label{tableau:mamba_hyperparametre}
\end{table}

\section{Séparation des données de test et d'entraînement}

\section{Entraînement}
\subsection{Fonction de perte}
\subsection{Propagation et Rétropropagation}
\subsection{Pré-entraînement}
\subsection{Prévention du surajustement}

\section{Évaluation}

\section{Logiciel}

\section{Résumé des modèles}
